\section{Exact constraint-conditioned deal inference}
\label{sec:inference}

Every agent in this lineage rests on an inference object that this section defines formally. The
method uses exact enumeration through a factorised dynamic program. The object it computes is the
uniform distribution over deals satisfying the acting player's hand information and the hard
constraints induced by the public history; \S\ref{sec:inference-prior} distinguishes that object
from a policy-aware posterior and describes the approximation the agent deploys.

\subsection{Notation}
\label{sec:inference-notation}

Table~\ref{tab:notation} fixes the symbols used here and in \S\ref{sec:agent}.

\begin{table}[!t]
\centering
\caption{Notation used in \S\ref{sec:inference} and \S\ref{sec:agent}.}
\label{tab:notation}
\tabfont
\begin{tabularx}{\linewidth}{@{}l L@{}}
\toprule
Symbol & Meaning \\
\midrule
\multicolumn{2}{@{}l}{\itshape The game and the hidden state} \\
$N=6$, $M=9$, $|H|=6$ & seats, half-suits, cards per half-suit; the deck is $M|H|=54$ cards \\
$s$ & the acting seat; $\mathcal{T}_s$ its two teammates, $\mathcal{O}_s$ its three opponents \\
$D$ & the hidden state: the initial deal, an assignment of all 54 cards to seats, nine each \\
$\mathcal{I}_s^t$ & seat $s$'s information set after $t$ public events \\
$\mathcal{C}_s^t$ & the hard constraints (C1)--(C5) seat $s$ has accumulated after $t$ events \\
$\mathcal{P}_t$ & a partial digest (rolling hash) of the public event stream (\S\ref{sec:agent-tiebreak}) \\
$\mathrm{own}_t(c)$, $H(c)$ & the seat holding card $c$; the half-suit containing $c$ \\
$\mathcal{H}_t$ & the half-suits not yet claimed \\
\midrule
\multicolumn{2}{@{}l}{\itshape Beliefs} \\
$\mathcal{M}_t(c)$ & the candidate set of card $c$: seats not yet excluded by a certificate \\
$U_t$, $n_p$ & the unresolved cards; the unresolved cards seat $p$ holds \\
$B_h$, $\mathcal{F}_h$ & the block of half-suit $h$; its assignments satisfying every certificate \\
$\sigma(\alpha)$, $T_h(\mathbf{v})$ & an assignment's per-seat count signature; the block count table \\
$F_k$, $G_k$ & the forward and backward partial convolutions of the count join \\
$q_s^t(D)$, $Z$ & the constraint-conditioned distribution of \eqref{eq:posterior}; the number of
  consistent deals \\
$\mu(c,p)$ & the exact marginal of $q_s^t$: the probability that seat $p$ holds card $c$ \\
$\hat\mu(c,p)$ & the deployed approximation to $\mu(c,p)$, with behavioural weighting \\
$\theta$, $\varphi$ & the two behavioural-weighting parameters of \eqref{eq:prior} \\
\midrule
\multicolumn{2}{@{}l}{\itshape Acting} \\
$\mathcal{A}_t$, $a=(c,q)$ & the legal ask set; an ask naming card $c$ and opponent $q$ \\
$p_a$ & $\hat\mu(c,q)$, the modelled hit probability of ask $a$ \\
$\phi(a)$, $w$, $u(a)$ & the ask feature vector, its fitted weights, and the resulting score \\
$\lambda$, $\beta$ & the linear and value scale parameters of \eqref{eq:askscore} \\
$M_H$, $q_H$ & opponent-held mass and residual ambiguity in half-suit $H$, in \eqref{eq:contest} \\
$\mathcal{E}_t$, $\varepsilon$ & the bit-for-bit tied candidate set and its tolerance \\
$\pi_{\mathrm{alloc}}$, $\pi_{\mathrm{team}}$ & under $q_s^t$: the probability a named allocation
  is exactly correct; that the team owns a half-suit outright \\
$\hat\pi_{\mathrm{alloc}}$, $\hat\pi_{\mathrm{team}}$ & the deployed approximations to those two
  quantities (\S\ref{sec:agent-declare}) \\
$\delta$, $\tau$ & the declaration value margin and the forced-endgame threshold \\
$\chi_t$, $\nu_t$, $K$ & the deduction-state digest, the no-progress counter, its threshold \\
$\Phi_t$ & the potential function of the termination argument in \S\ref{sec:agent-stall} \\
$D_{\mathrm{det}}$, $w_d$ & the number of sampled worlds and their importance weights \\
$\kappa$, $\eta$ & the search's deviation threshold and its blueprint blend, $\eta = 0$ here \\
\bottomrule
\end{tabularx}
\end{table}

\subsection{Engine architecture and information access}
\label{sec:inference-engine}

The engine is a single-header C++20 core. A hand is one 64-bit word over the 54 card indices
$c = 6h + i$ for half-suit $h$ and offset $i$, so hand arithmetic is bit arithmetic. Random streams
are counter-based and keyed by the triple (seed, seat, purpose). This gives the property the
duplicate-block design of \S\ref{sec:evaluation-duplicate} requires: a deal is reproducible from
its seed, and the play of a deal between deterministic agents reproduces from the seed and the
agents' configurations. Complete reproducibility of a scored run is a stronger property, and
\S\ref{sec:agent-coordination} and Appendix~\ref{app:corrections-self} record where searching
agents fall short of it.

An acting policy never receives the game state. It receives its own knowledge object, its own hand
and the public event stream. There is no reference from a policy to the deal, to another seat's hand
or to the referee. This is an architectural barrier rather than a convention, and it is what allows
the side-channel audit of \S\ref{sec:evaluation-side} to be read as a test of the policy: a policy
that cannot reach the deal cannot read it, so a replay failure indicates dependence on information
or state not represented in the replay construction. Subsequent source inspection localised the
observed residual to persistent rollout state (\S\ref{sec:agent-coordination}), although the
specific carrier field and its scored win-rate effect were not isolated.

\subsection{Observer-specific deal constraints}
\label{sec:inference-constraints}

Because a card changes hands only through an observed ask, the public record determines
$\mathrm{own}_t(\cdot)$ given $D$. The only hidden object is therefore $D$ itself. Write
$\mathcal{D}$ for the set of all deals; from one seat's viewpoint at deal time
$|\mathcal{D}| = 45!/(9!)^5 \approx 1.9 \times 10^{28}$.

Each public event imposes a \emph{hard} constraint, meaning a logical restriction on $\mathcal{D}$
rather than a probabilistic tilt. The engine maintains them as a per-card candidate set
$\mathcal{M}_t(c)$, initialised to all six seats. For every card that remains unresolved the set
only shrinks, $\mathcal{M}_{t+1}(c) \subseteq \mathcal{M}_t(c)$; a card whose holder becomes public
through a transfer leaves the unresolved set $U_t$ with its holder recorded, so the monotonicity
statement is about unresolved cards:

\begin{description}
\item[(C1) Own hand.] For each $c$ in seat $s$'s hand, $\mathcal{M}(c) \leftarrow \{s\}$; for every
  other card, $s$ is removed from $\mathcal{M}(c)$.
\item[(C2) Card counts.] Every seat's current hand size is public, so with the transfer record each
  seat's count of unresolved cards, $n_p$, is known. This is the only constraint that couples
  different half-suits.
\item[(C3) Transfers.] A successful ask makes $\mathrm{own}(c)$ public from then on, so
  $\mathcal{M}(c)$ collapses to a singleton.
\item[(C4) Failed asks.] If $q$ was asked for $c$ at time $r$ and did not hold it, then $q \notin
  \mathcal{M}(c)$ unless $q$ has acquired $c$ since, which is public.
\item[(C5) Ask legality.] An asker of $c$ held another card of $H(c)$ at that moment and did not hold
  $c$. The first is a disjunction over the other five cards of the half-suit and is the one
  constraint that is not a simple exclusion.
\end{description}

Write $\mathcal{C}_s^t$ for the constraints seat $s$ has accumulated. The object this section
computes is the uniform distribution over the deals that satisfy them:
\begin{equation}
  q_s^t(D) \;=\; \frac{\ind\!\left[D \models \mathcal{C}_s^t\right]}{Z},
  \qquad
  Z \;=\; \bigl|\{D' \in \mathcal{D} : D' \models \mathcal{C}_s^t\}\bigr|.
  \label{eq:posterior}
\end{equation}
This distribution is exact for the stated constraint model. It is not, in general, the posterior
$\Pr[D \mid \mathcal{I}_s^t]$: a uniform prior over initial deals does not imply a uniform posterior
after strategically chosen actions, because different feasible deals can make the observed action
sequence differently likely under the players' policies. A policy-aware posterior would weight each
feasible deal by the likelihood of the observed actions under a specified behavioural model, which
is the computation the transcript-inversion adversary of \S\ref{sec:results-whitebox} approximates
and \eqref{eq:coord} states. The modelling content of \eqref{eq:posterior} is entirely in which
constraints are imposed; nothing in it refers to how anyone plays. \S\ref{sec:inference-prior}
describes the behavioural weighting the deployed approximation adds.

\subsection{Half-suit factorisation and count dynamic programming}
\label{sec:inference-factor}

Computing $Z$ by enumerating $\mathcal{D}$ is infeasible, and a constraint system of this shape
generally induces a permanent-like count. For such counts the literature supplies exact formulae
that are exponential in the general case, inclusion--exclusion \cite{ryser} and the sign-vector
formula \cite{glynn}, and approximation algorithms, Markov chains \cite{jsv}, matrix scaling
\cite{lsw} and sequential importance sampling \cite{chen-sis}.

Fish avoids that, because every constraint except (C2) is confined to a single half-suit. Constraints
(C3), (C4) and (C5) all speak about one card or one half-suit; only the per-seat capacities $n_p$
couple half-suits, and they couple them only through counts.

Partition the unresolved cards by half-suit, $U_t = \bigcup_{h} B_h$, where $B_h$ is the
\emph{block} of half-suit $h$. For a block, let an \emph{assignment} be a map $\alpha: B_h \to
\{0,\dots,5\}$ with $\alpha(c) \in \mathcal{M}(c)$, and let $\mathcal{F}_h$ be the assignments that
additionally satisfy every (C5) disjunction touching $h$. Write the \emph{count signature} of an
assignment as the vector
\[
  \sigma(\alpha) \;=\; \bigl(|\alpha^{-1}(0)|,\;\dots,\;|\alpha^{-1}(5)|\bigr) \in \mathbb{Z}_{\ge 0}^{6}.
\]

\paragraph{Block enumeration.} For each half-suit independently, enumerate $\mathcal{F}_h$
and accumulate a table indexed by count signature,
\begin{equation}
  T_h(\mathbf{v}) \;=\; \bigl|\{\alpha \in \mathcal{F}_h : \sigma(\alpha) = \mathbf{v}\}\bigr|,
  \qquad
  T_h^{c \to p}(\mathbf{v}) \;=\; \bigl|\{\alpha \in \mathcal{F}_h : \sigma(\alpha) = \mathbf{v},\,
  \alpha(c) = p\}\bigr|.
  \label{eq:blocktable}
\end{equation}
A half-suit holds six cards and resolved ones are removed, so $|B_h| \le 6$ and the enumeration is
small.

\paragraph{Count convolution.} Constraint (C2) says the signatures must sum to the vector of
per-seat capacities $\mathbf{n} = (n_0,\dots,n_5)$. Order the blocks arbitrarily as $h_1, \dots,
h_M$ and define the partial convolution. A half-suit carrying no live (C5) disjunction factorises
further into one independent unit per card, so the implementation runs one layer per unit rather
than one per half-suit; the recursion below is the same object either way.
\begin{equation}
  F_0(\mathbf{v}) = \ind[\mathbf{v} = \mathbf{0}], \qquad
  F_k(\mathbf{v}) \;=\; \sum_{\mathbf{u}} F_{k-1}(\mathbf{v}-\mathbf{u})\, T_{h_k}(\mathbf{u}),
  \label{eq:dp}
\end{equation}
with every table taken as zero outside the box $\mathbf{0} \le \mathbf{v} \le \mathbf{n}$
(componentwise), so that
\begin{equation}
  Z \;=\; F_M(\mathbf{n}).
  \label{eq:Z}
\end{equation}
By induction on $k$, $F_k(\mathbf{v})$ counts the joint assignments of the first $k$ blocks that
satisfy their own certificates and have total signature $\mathbf{v}$; since the blocks partition
$U_t$ and (C2) is the only constraint that couples them, $F_M(\mathbf{n})$ counts exactly the
assignments satisfying (C1)--(C5), which is $Z$. Running the same recursion backwards, with
$G_{M+1}(\mathbf{v}) = \ind[\mathbf{v} = \mathbf{0}]$ and $G_k(\mathbf{v}) = \sum_{\mathbf{u}}
T_{h_k}(\mathbf{u})\, G_{k+1}(\mathbf{v}-\mathbf{u})$, gives suffix tables $G_k$, and the marginals
follow by substituting one block's card-resolved table for its plain one:
\begin{equation}
  \mu(c,p) \;=\; \frac{1}{Z}\sum_{\mathbf{u},\,\mathbf{v}}
  F_{k-1}(\mathbf{v})\; T_{h_k}^{c \to p}(\mathbf{u})\; G_{k+1}(\mathbf{n}-\mathbf{v}-\mathbf{u}),
  \qquad c \in B_{h_k}.
  \label{eq:marg}
\end{equation}

Three quantities the policy needs are \emph{defined} by this machinery, and
\S\ref{sec:inference-prior} records that the deployed configuration substitutes a cheaper
approximation for each of them. The per-card marginals are \eqref{eq:marg}. The probability under
$q_s^t$ that the acting team owns half-suit $h$ outright is the total weight of the block
assignments whose image lies in the team,
\begin{equation}
  \pi_{\mathrm{team}}(h) \;=\; \frac{1}{Z}
  \sum_{\substack{\alpha \in \mathcal{F}_h\\ \alpha(B_h) \subseteq \mathcal{T}_s \cup \{s\}}}
  \;\sum_{\mathbf{v}} F_{k-1}(\mathbf{v})\; G_{k+1}\bigl(\mathbf{n}-\mathbf{v}-\sigma(\alpha)\bigr),
  \qquad h = h_k.
  \label{eq:team}
\end{equation}
And the probability of a fully named allocation, which is exactly what a declaration asserts, is
\begin{equation}
  \pi_{\mathrm{alloc}}(h, \alpha^\star) \;=\;
  \frac{\ind\!\left[\alpha^\star \in \mathcal{F}_h\right]}{Z}
  \sum_{\mathbf{u}, \mathbf{v}} F_{k-1}(\mathbf{v})\,
  \ind\!\left[\sigma(\alpha^\star) = \mathbf{u}\right]
  G_{k+1}(\mathbf{n}-\mathbf{v}-\mathbf{u}),
  \label{eq:alloc}
\end{equation}
for the single assignment $\alpha^\star$ the declaration names. The leading indicator is not
cosmetic: an assignment with a surviving count signature need not itself satisfy the half-suit's
certificates, so feasibility is checked explicitly. No sampling is involved in any of the three.

\paragraph{Cost.} The signature space is bounded by the capacities, each block table has at most
$6^{|B_h|}$ entries before pruning, and the join is $M$ convolutions over that space. An exact
rebuild was timed at a mean of \vsixCountCost{}~ms in the \vsix{} cycle's measurement, taken over
the exact-program builds at contested ask decisions in 150 mirror self-play games on the desktop
described in Appendix~\ref{app:repro}; the cost is strongly state-dependent and the figure is a
mixed-regime mean. The counts are exact as a formulation; the implementation evaluates the block
tables in single precision and the convolutions, $Z$ and the marginals in double precision with no
rescaling, so the values are floating-point evaluations of exact integer quantities, and $Z$ at deal
time exceeds the range in which a double represents integers exactly. The verified calculation is the constraint-conditioned count and
its marginals; it is not a policy-aware posterior, and the project's decision to spend its cycles on
policy rests on the former being cheap rather than on inference being solved in the latter sense.

\subsection{Approximate inference and behavioural weighting}
\label{sec:inference-prior}

Equation~\eqref{eq:posterior} is exact \emph{for the constraint model}: it is the uniform
distribution over constraint-consistent deals. It contains no model of behaviour, because the
construction above consumes the deduction state alone and has no parameter through which an
opponent's tendencies could enter.

The path the agent deploys adds two such parameters and gives up exactness for them. The
history-dependent weighting it introduces is called the \emph{policy prior} in the code and in
earlier reports; this paper calls it a behavioural weighting model, since it is a reweighting of the
constraint-consistent deals by the observed ask history rather than a prior in the Bayesian sense.
It replaces
the count join by iterative proportional fitting, better known here as Sinkhorn scaling, on the
$|U_t| \times 6$ matrix of admissible card--seat pairs. Initialise
\begin{equation}
  \hat\mu^{(0)}(c,p) \;=\; \ind\!\left[p \in \mathcal{M}(c)\right] \cdot
  \exp\Bigl(\operatorname{clip}\bigl(\theta\, a_{p,H(c)} - \varphi\, (A_p - a_{p,H(c)}),\; \pm 2.6\bigr)\Bigr),
  \label{eq:prior}
\end{equation}
where $a_{p,h}$ is the number of asks seat $p$ has made in half-suit $h$ and $A_p$ its total asks, and
then alternate the two normalisations
\begin{equation}
  \hat\mu(c,\cdot) \leftarrow \frac{\hat\mu(c,\cdot)}{\sum_{p}\hat\mu(c,p)},
  \qquad
  \hat\mu(\cdot,p) \leftarrow \hat\mu(\cdot,p)\cdot\frac{n_p}{\sum_{c}\hat\mu(c,p)},
  \label{eq:sinkhorn}
\end{equation}
to a fixed iteration count. The row step makes each card's distribution a probability; the column
step enforces the capacity constraint (C2) that the exact join enforces combinatorially.

Between outer sweeps the routine also applies an independence conditioning step for each live (C5)
disjunction: for a certificate saying that seat $j$ held \emph{some} card of a half-suit,
each candidate card $c'$ of that half-suit is rescaled to
$\hat\mu(c',j) \big/ \bigl(1 - \prod_{c''}(1 - \hat\mu(c'',j))\bigr)$, the product running over the
half-suit's cards $c''$.
So (C5) is neither enforced exactly nor discarded in the deployed path: it is imposed approximately,
under an independence assumption the exact enumeration of \S\ref{sec:inference-factor} does not need.
The deployed parameters are $\theta = \vsevenPriorTheta{}$ and $\varphi = \vsevenPriorPhi{}$.

Equation~\eqref{eq:prior} is the entire behavioural content of the agent's beliefs: a seat that has
asked repeatedly in a half-suit is treated as more likely to hold cards of it, because asking
requires already holding one. It is a soft version of the (C5) certificate, in the spirit of the Skat
inference line \cite{solinas-supervised,rebstock-inference,buro-kermit} but far simpler, with two
parameters and no learned opponent model.

\S\ref{sec:background-vsix} reports the model comparison: scored as predictors of card location on
identical states, $\hat\mu$ has lower log loss and a higher argmax hit rate than $\mu$. In this
implementation, exactness for the constraint model and behavioural weighting are in different
objects: the block program carries no behavioural parameter, and the deployed paths add the
weighting either as the initialiser of the Sinkhorn fit or, in the search, as an importance weight
on exact samples (\S\ref{sec:agent-search}). In the comparison, the weighting contributed more
predictive accuracy than exactness did. That is why the
paper treats exact constraint-based inference as an assumption to test rather than a property to
assume; \S\ref{sec:development-alloc} reports a second comparison of the same kind.

\subsection{Implementation verification}
\label{sec:inference-verify}

Three checks run over the engine, and all three ran as preconditions of the evaluation in
\S\ref{sec:results}.

\begin{description}
\item[Soundness audit.] Every seat's candidate set is checked against the true deal at every
  decision: $\mathrm{own}_t(c) \in \mathcal{M}_t(c)$ must hold always, since no agent's deduction may
  exclude the truth. The verification run for this evaluation recorded
  \vsevenVerifyViolations{} violations in \vsevenVerifyChecks{} checks, and the preceding cycle
  \vsixAuditViolations{} in \vsixAuditChecks{}.
\item[Cross-engine agreement.] The block dynamic program, an independent card-counting path and an
  exact rejection sampler must agree on \eqref{eq:marg}. The self-test reports
  \vsevenSelftestChecks{} checks passing, a check being one (card, observer, event) row. Its
  tolerances are $10^{-9}$ for block-versus-card-program agreement, which is checked on states with
  no live (C5) disjunction and holds there at the limit of double precision, and $0.05$ for
  block-versus-rejection-sampling agreement on states with live disjunctions, which is Monte Carlo
  precision at the sample sizes the test draws. A separate brute-force oracle checks $Z$, the
  marginals, and the team and allocation probabilities against exhaustive enumeration on small
  states to $10^{-9}$.
\item[Identity control.] The current policy, carrying its predecessor's parameter vector with every
  new mechanism disabled, must produce a transcript identical to its predecessor's under md5. It does
  (\vsixIdentity{}). This control checks whether refactoring changes behaviour before the new
  components are enabled, so that a measured ablation difference is not an artefact of refactoring.
\end{description}

Two defects these checks found are recorded here because both changed previously published numbers,
and both survived for the same reason: they were invisible to the path the agent actually deploys.
The block dynamic program parked every instance's tables in a shared per-thread pool, so a second
construction on the same thread silently repointed the first instance's tables, giving
\numAliasMismatch{} mismatches in \numAliasChecks{} checks; any earlier number computed with the
exact belief in a configuration where more than one seat held such an object is affected, and the
deployed approximate-belief numbers are not. Separately, an audit that verifies the declaration
pre-gate never rejects a correct declaration had been parsed only in an older branch and so had never
run; once enabled it reports \vsixGateFalseNeg{} false negatives over \vsixGateOpportunities{}
declaration opportunities. Appendix~\ref{app:corrections} carries the full register.
